%!TEX root = ../main.tex
\chapter{Training Methodology}
\label{ch:training}

\section{Two-Stage Strategy and Rationale}
\label{sec:training-strategy}

\begin{figure}[htbp]
  \centering
  \linespread{1}\selectfont
  \begin{tikzpicture}[x=1cm, y=1cm, blk/.append style={minimum width=34mm}]
    \node[ext, minimum width=34mm] (corpus) at (0,0)    {recorded corpus};
    \node[lyr, minimum width=34mm] (sup)    at (0,-1.35) {imitation\\\footnotesize cap-scaled MSE};
    \node[blk] (w)      at (0,-2.70) {weights};
    \node[lyr, minimum width=34mm] (rl)     at (0,-4.05) {reinforcement\\\footnotesize recurrent PPO};
    \node[blk] (pol)    at (0,-5.40) {policy};

    \draw[flow] (corpus) -- (sup);
    \draw[flow] (sup) -- (w);
    \draw[flow] (w) -- (rl);
    \draw[flow] (rl) -- (pol);

    \node[ext, minimum width=30mm] (fleet) at (-4.4,-4.05)
      {fleet of parallel slices\\\footnotesize the pipeline, $N$ times};
    \draw[thin flow] (fleet) -- (rl);

    \node[note, anchor=west, text width=36mm] at (2.1,-1.35)
      {learns on the states the\\demonstrator visited};
    \node[note, anchor=west, text width=36mm] at (2.1,-2.70)
      {already deployable\\on the robot};
    \node[note, anchor=west, text width=36mm] at (2.1,-4.05)
      {learns on the states it\\generates itself};
  \end{tikzpicture}
  \caption{The two stages of Section~\ref{sec:training-strategy}. The first
           produces both the weights deployed on the robot and the starting
           point of the second; the second fine-tunes them on the distribution the
           policy itself generates, which is where the behaviour the local map
           should determine finally has a consequence.}
  \label{fig:two-stage}
\end{figure}

The policy described in Chapter~\ref{ch:lnn} is trained twice
(Figure~\ref{fig:two-stage}), and the two stages answer two limitations that
compensate for one another.

Imitation learns, from recorded routes, the command the pilot gave from the
observation the robot had. It is inexpensive --- the data already exist,
produced by the pipeline during normal use of the system --- and yields a
sensible starting behaviour. It carries, however, the structural defect recalled
in Section~\ref{sec:rel-imitation-rl}: the policy is evaluated on the
distribution of the states the demonstrator visited, whereas in execution it
visits those produced by its own actions. The defect is aggravated by the nature
of the demonstration: the pilot, human or automatic, plans its route on the true
geometry of the world, so that its command depends on the local map only
indirectly. Imitation teaches the reproduction of a behaviour, not the reaction
to what is observed.

Reinforcement learning does not have that defect, since it learns on the
distribution it generates itself and judges actions by the consequences they
produce. It requires, however, a number of interactions a real robot cannot
provide, each of which may destroy the vehicle. It is therefore carried out in
simulation, on a fleet of parallel environments.

The strategy adopted is the composition of the two: the first stage produces the
weights the robot flies with and the starting point of the second; the second
fine-tunes them on the correct distribution, where the behaviour the local map
ought to determine --- avoiding an obstacle --- finally has a consequence.

\paragraph{What the two stages share.}
The two procedures are implemented separately and share exactly three things:
the architecture table of Section~\ref{sec:module-implementation}, the way a
point cloud is packed for the encoder, and the thermal guard that suspends
training when the temperature of the machine exceeds a threshold. The sharing is
deliberately minimal and deliberately no greater: how the weights were fed is
part of what the weights mean, and two implementations of the same packing would
diverge without raising any error.

Both stages are moreover validated on simpler tasks before being applied to the
problem --- a synthetic task for the first, a standard control environment for
the second --- according to the protocol described in
Section~\ref{sec:exp-protocol}.

\section{Supervised Stage}
\label{sec:training-supervised}

\subsection{Corpus Acquisition and Recording Protocol}
\label{sec:sl-corpus}

The corpus is recorded through the pipeline itself, by the node described in
Section~\ref{sec:observation-recording}: every sample pairs the observation
available to the vehicle with the command that followed from it. A property no
other arrangement could guarantee follows from this: the distribution on which
the network is trained is produced by the same code that will feed it on board.

One episode corresponds to one directory. A mission may record a flight whole or
in two segments --- the outward and the return leg --- each of which becomes an
episode of its own while sharing with the other the world, the starting position
and the goal. It is a circumstance the data split must be aware of, and
Section~\ref{sec:sl-split} returns to it.

Every episode carries a metadata file describing it: the robot, the world, the
voxel pitch, the list of the archives that compose it, and a statistical summary
of the recorded commands. The summary is accumulated \emph{while} writing rather
than derived afterwards, because deriving it afterwards would mean reopening
every archive of the corpus: an operation of seconds over ten gigabytes and
impracticable over a terabyte, which would moreover have to be repeated at every
change of seed.

\subsection{Dataset Format: Sharded Episodes and Manifest}
\label{sec:sl-dataset}

Episodes are stored as sequential archives, read as a stream rather than loaded
into memory. The format imposes a property that conditions the structure of the
sample: since the archives are read in shuffled order, every sample must stand
on its own. The state block is therefore repeated in each sample although it is
constant within a recording --- a few dozen bytes against the tens of kilobytes
occupied by the voxels --- rather than being stored once per episode.

Building the \emph{manifest} of the corpus, that is the list of the episodes
with their properties, reads the metadata files exclusively and never opens an
archive. It is this separation that makes it possible to recompute a data split
in negligible time.

\subsection{Stratified Split and Episode Balancing}
\label{sec:sl-split}

Splitting per \emph{sample} would produce a meaningless evaluation: adjacent
frames of the same flight are nearly identical, and putting some of them in
training and the others in validation means measuring training and calling it
validation.

Splitting per \emph{episode} repeats the same error one level higher, for the
reason set out in Section~\ref{sec:sl-corpus}: the two segments of a flight
recorded in two parts share world, start and goal, and would end up on opposite
sides of the split. The unit adopted is therefore the \emph{run}, declared in
the metadata, of which a flight recorded whole is the particular case with a
single segment.

Runs are moreover stratified on what changes the distribution of the data: the
robot model, since different platforms command different things; the world,
since the statistics of the obstacles change; the voxel pitch, since the
quantisation changes and with it the extent of the grid
(Section~\ref{sec:observation-cloud}); and the set of command components the
episode actually exercises, since an episode in which the robot never turns
teaches nothing about turning.

Two provisions complete the protocol. The split is governed by a seed of its
own, distinct from the one that initialises the weights and the shuffling:
repeating a session under a different seed therefore measures what that session
is worth and not what a different split is worth. And the test partition is read
once, at the end, since consulting it between sessions would turn it into a
second validation partition.

\subsection{Sequence Chunking and Hidden-State Propagation}
\label{sec:sl-sequences}

A recurrent network is trained on windows of consecutive samples. The length of
the window fixes the horizon over which the gradient propagates backwards; the
step between the start of one window and the start of the next determines
instead how many windows are obtained from the same recordings, and a step
shorter than the length makes them overlap, which is the cheapest answer
available to a small corpus.

A window is cut wherever the interval between two samples exceeds a threshold: a
sequence that straddles a prolonged silence is not a sequence. The first step of
every window is moreover given the median interval of the window itself rather
than a null value, because a null interval pins the gate of the CfC cell at
exactly $0.5$ --- the degeneration discussed in Section~\ref{sec:cell-cfc} ---
precisely at the step where the state has just been cleared.

The recurrent state is governed by the mechanism of
Section~\ref{sec:module-implementation}: cleared before every window, because
each is an independent sample of the distribution of trajectories and carrying
into it the memory of the previous one would mean training on a history the
sample does not have; and detached from the graph after every optimiser step,
because otherwise the graph would grow from one mini-batch to the next until it
covered the whole epoch.

\paragraph{The composition of the batch.}
The sparse convolutional operator accepts a single grid shape per call, so that
a batch may contain only samples of the same grid. The most immediate
realisation --- one episode per batch --- satisfies the constraint but imposes a
stronger one than necessary: with overlapping windows a batch then holds few
distinct scenes, all of the same recording, and the encoder has nothing to
contrast within it. The loader therefore keeps several archives open at once and
fills one pending batch per grid drawing from all of them, so that the effective
constraint is the grid and not the episode. Validation and test keep one episode
per batch regardless, so that their figures remain comparable across sessions.

\subsection{Data Augmentation}
\label{sec:sl-augmentation}

The only transformation applied to the data is the left--right reflection of a
training window, through the longitudinal plane of the body, before it is
batched. The reflected flight is as valid as the recorded one for a vehicle
symmetric about its own longitudinal axis, and it doubles the scenes the encoder
observes on each side.

The reflection is correct only if applied to \emph{everything} that has a side,
and this is the delicate point of the transformation: the lateral index of the
cloud about the centre of the grid, the lateral component of the pose, of the
goal and of the centre of mass, the lateral velocity and command, and the axial
quantities --- the roll and yaw rates, and the corresponding components of both
quaternions. The targets of the auxiliary supervision of
Section~\ref{sec:sl-geometry} swap sector $k$ for sector $-k$. A window is
reflected whole and never step by step, and validation never is.

\subsection{Observation Normaliser Fitting}
\label{sec:sl-normaliser}

The affine transform described in Section~\ref{sec:aux-modules} is estimated
from the data before training, over a bounded number of batches, and stored
inside the weights. It follows that inference on board applies it without
knowing: there is no normalisation step external to the network that a node
might forget to perform.

Features of zero variance require explicit treatment. Since the normalisation
divides by the standard deviation bounded below by a minimum value, a feature
constant over the corpus would be divided by that minimum, and the day it ceased
to be constant --- a tilted goal, a limit changed in the configuration --- it
would enter the network amplified by six orders of magnitude. The deviation of
such features is therefore raised to one. The correction is neutral on constant
data, since a deviation from the mean that is identically zero remains zero
whatever the divisor.

\subsection{Cap-Scaled Mean Squared Error Loss}
\label{sec:sl-loss}

The mean squared error computed on the raw command is dominated by the component
that carries the most energy, and the six components of a velocity command
differ from one another by orders of magnitude in variance: the gradient
reaching the small ones is negligible, and they are in effect never learnt.

The customary correction --- dividing by the per-dimension standard deviation
measured on the corpus --- is not applicable here, for two reasons. Some
components are identically zero over the whole corpus, so their scale would be
an arbitrary value of convenience. And every scale would shift as new recordings
arrived, silently reweighting the loss function between sessions meant to be
comparable.

The loss adopted divides instead by the \emph{velocity limits declared by the
robot}, which are already present in the state block of the observation
(Section~\ref{sec:observation-vector}); its formal definition is given in
Section~\ref{sec:metrics-loss}. Three properties follow. The error becomes a
fraction of what that robot can do, so that platforms with very different limits
are directly comparable. The scale is a property of the platform and not of the
corpus, and does not move when data is added. And a heterogeneous fleet requires
no special case, since every sample carries its own limits.

\subsection{Auxiliary Geometry Supervision}
\label{sec:sl-geometry}

For the reason set out in Section~\ref{sec:geometry-head}, imitation alone
provides the encoder with a weak signal. The auxiliary supervision adds a second
objective that requires no labels.

The targets are measured once per corpus, before training, on the cloud of each
frame: the distance to the nearest obstacle in each of eight azimuth sectors and
the height above the ground. The measurement is performed with the same
functions the perception bridge uses on board --- the same estimate of the
ground and the same exclusions of the hull of the vehicle --- so that the
encoder is taught exactly the quantity the reward of the reinforcement stage
prices (Section~\ref{sec:rl-reward}). The result is written beside each episode.

The overall loss is the sum of the imitation loss and the auxiliary one,
weighted by a coefficient. Two properties of the protocol must be stated because
they make sessions with and without the auxiliary term comparable. The first is
that the losses recorded as training and validation loss remain the imitation
loss alone, so that two sessions are ranked on the same quantity. The second is
that the output of the encoder is computed once per batch and handed to the
network, rather than being recomputed for the auxiliary head: the sparse
convolution, which is the expensive part of the forward pass, is not performed
twice.

\subsection{Optimisation and Checkpoint Selection}
\label{sec:sl-optimisation}

The optimiser is AdamW, with a weight penalty sized as a regulariser rather than
as a constraint. The learning rate decays linearly towards a fraction of its
initial value rather than towards zero, so that the last epochs keep learning
instead of merely not getting worse; the gradient norm is clipped at every step.

The horizon of the decay requires knowing the number of steps per epoch, which a
streaming loader does not declare: it is supplied as an initial estimate and
measured for real on the first pass, after which the schedule is retargeted onto
the observed value.

\textbf{No automatic early stopping is provided.} Every epoch produces a
checkpoint, and the epoch that minimises the validation loss is flagged:
selecting the checkpoint is part of the experimental protocol
(Section~\ref{sec:exp-protocol}) and not a decision taken by the training. The
choice is deliberate --- stopping automatically would mean fixing the criterion
once and for all, whereas the criterion is what Chapter~\ref{ch:results} must be
able to discuss. A checkpoint may accordingly be reloaded in two ways: with its
weights alone, which is what starting a new session from them requires and is
the case of the reinforcement stage, or in full, which also restores the state
of the optimiser and the position in the schedule and is what continuing an
interrupted session requires.

Two seeds govern reproducibility: one initialises the weights and the shuffling,
the other the data split alone (Section~\ref{sec:sl-split}). A thermal guard,
finally, suspends training when the temperature of the machine exceeds a
threshold and resumes it only once it has fallen by a margin, so that a machine
at the limit does not alternate suspension and resumption at every check.

\subsection{Output Calibration Fitting}
\label{sec:sl-calibration}

The recalibration of the output, motivated in
Section~\ref{sec:output-calibration}, is estimated according to a precise
protocol, which belongs to this chapter while its purpose belongs to the
previous one.

Validation is run with the calibration stage at the identity, and the
calibration is estimated \emph{from those same predictions}, afterwards. The
ordering is not arbitrary: a model is already calibrated with respect to the
data on which it was optimised, so that an estimate made on the training set
would read a gain close to one and would under-correct everywhere else.

The estimate is repeated at every epoch, so that each checkpoint carries the
calibration belonging to its own weights. A gain exceeding a declared maximum is
clipped and reported: if the spread of the prediction is a small fraction of
that of the command, the least-squares line is extrapolating.

\section{Reinforcement Stage}
\label{sec:training-rl}

\subsection{Problem Formulation}
\label{sec:rl-pomdp}

The task is formulated as a partially observable Markov decision process, and it
is worth identifying its components on the concrete problem rather than in the
abstract.

The \emph{state} comprises the configuration of the vehicle --- pose, velocity,
attitude --- and that of the environment surrounding it; the agent never has
access to it. What it receives is the \emph{observation} defined in
Section~\ref{sec:observation-space}, produced by the same pipeline that produces
it on board, and which is a partial and noisy function of that state. The
\emph{action} is the six-component velocity command, continuous and expressed in
the physical units of the platform. The \emph{transition} is determined by the
dynamics of the simulated vehicle and by the latency of the chain that produces
the observation: between the measurement on which the agent decides and the
effect of the command lies the traversal time of the pipeline, which is part of
the problem rather than a defect to be compensated for. The \emph{reward} is
discussed in Section~\ref{sec:rl-reward}, the episode horizon in
Section~\ref{sec:rl-episodes} and the discount factor in
Section~\ref{sec:rl-presets}.

\paragraph{What the partial observability consists of.}
It is not a formality: on this problem it has three distinct origins, and all
three have consequences for the learnt behaviour. The local map represents what
the sensors have observed and not yet forgotten
(Section~\ref{sec:observation-map}), so that a region never framed is
indistinguishable from a free one. The pose is an estimate subject to drift
rather than the true position, and the goal is expressed relative to it. A single
frame, finally, says nothing about what was happening: whether the vehicle was
accelerating, whether the obstacle ahead is approaching or receding, how long it
has been standing still.

\paragraph{The policy is a function of the history.}
It follows that, in principle, the optimal policy is a function not of the last
observation but of the sequence preceding it; how much of that history a policy
actually needs on a given task is an empirical question, which the stateless
control answers in Chapter~\ref{ch:results}. The two customary treatments are to stack a window
of observations into a single vector, or to maintain a state updated at every
step; the architecture of Chapter~\ref{ch:lnn} adopts the second, so that the
policy is $\pi_\theta(a_t \mid o_t, x_t)$, with $x_t$ the recurrent state
summarising the history. It is this choice that makes a recurrent version of the
algorithm necessary, and Section~\ref{sec:rl-buffer} describes its consequences.

\paragraph{The steps are not equally spaced in time.}
The formulation assumes that a step is a decision, but the duration of a step is
not constant: it depends on the rate at which the pipeline delivers the
observation (Section~\ref{sec:observation-recording}). The interval actually
elapsed is therefore part of what the policy reads, and it is the quantity with
which the cell integrates its own dynamics (Section~\ref{sec:cell-time}). The
same quantity is used by the environment to size the episode, which is defined
in steps but measured over a time (Section~\ref{sec:rl-episodes}).

\paragraph{The platform is not part of the hidden state.}
A last property distinguishes this formulation from the customary one of a
learnt navigation task, and it is the reason why a single policy can be trained
on a heterogeneous fleet. The identity of the body is not a latent variable the
agent must infer from the behaviour of the vehicle: it is declared in the
observation, in the kinematic limits and in the model identifier of
Section~\ref{sec:observation-vector}. Changing platform is therefore a change of
observation and not of problem --- the action space, the reward function and the
objective remain the same --- and it is what allows the experience of different
robots to be gathered into a single training.

\paragraph{The objective.}
The policy maximises the expected discounted return over the episodes generated
by the fleet. During training it is stochastic, since exploration is obtained by
sampling from a distribution whose width is itself trained
(Section~\ref{sec:rl-actions}); that width, however, belongs to the algorithm
and not to the network, so that the policy transferred on board is the mean of
that distribution, that is deterministic.

\subsection{Recurrent PPO: Objective and Clipping}
\label{sec:rppo-gae}

Policy-gradient methods improve the policy by raising the probability of the
actions that turned out better than expected. In the naive form this is valid
only for the data collected by the current policy: every batch can be used
exactly once, and a step too large destroys the policy without there being a
loss function to fall back on.

Proximal policy optimisation \cite{schulman2017ppo} buys reuse and safety by
keeping the updated policy \emph{close} to the one that collected the data.
Closeness is measured by the probability ratio of the action actually taken,
\begin{equation}
  r_t(\theta) = \frac{\pi_\theta(a_t \mid s_t)}{\pi_{\theta_{\text{old}}}(a_t \mid s_t)},
  \label{eq:ppo-ratio}
\end{equation}
and the objective is clipped so that drifting beyond a threshold ceases to be
profitable:
\begin{equation}
  L_{\text{policy}} = -\,\mathbb{E}\!\left[\min\!\left(r_t A_t,\;
     \operatorname{clip}(r_t, 1-\epsilon, 1+\epsilon)\, A_t\right)\right].
  \label{eq:ppo-clip}
\end{equation}

Equation~\eqref{eq:ppo-clip} reads as a hinge. If an action turned out better
than expected, raising its probability improves the objective until the ratio
reaches $1+\epsilon$, beyond which the clipped branch takes over and the
gradient vanishes: the update stops of its own accord. Symmetrically for an
action worse than expected.

An entropy bonus is subtracted from the objective, preventing the policy from
collapsing prematurely onto a single action. The critic is trained by regression
onto the returns, with the option of clipping its update as well with respect to
the value predicted during collection, so that a single mini-batch does not
rewrite the value function. Actor and critic have separate optimisers and
separate backward passes, and each receives its own clipping of the gradient
norm.

\subsection{Generalised Advantage Estimation and Truncation Bootstrapping}
\label{sec:rl-gae}

The advantage $A_t$ could be estimated as the difference between return and
value, but the return is a sum of many noisy rewards, while the one-step
estimate is biased by the critic's error. Generalised advantage estimation
\cite{schulman2016gae} interpolates between the two:
\begin{align}
  \delta_t &= r_t + \gamma\, V(s_{t+1})\,(1 - d_t) - V(s_t), \\
  A_t      &= \delta_t + \gamma\,\lambda\,(1 - d_t)\, A_{t+1},
  \label{eq:gae}
\end{align}
computed backwards over the rollout, where $d_t$ marks the end of an episode and
the factor $(1-d_t)$ prevents the trace from crossing its boundary.

\paragraph{Truncation and termination are not the same thing.}
An episode that has reached the time limit is not finished: the state in which
it stopped has a value, and treating it as zero introduces a systematic bias in
the critic precisely where episodes almost always end. The reward of a truncated
step therefore receives, during collection, the addition of $\gamma\,V$ computed
on the final observation --- which at that moment still exists, whereas the
automatic reset has already replaced it in the one returned. The advantage
estimate then uses the end-of-episode signal for both cases and does not repeat
the addition, which would count the value twice. An environment that truncates
without providing the final observation raises an error rather than silently
reintroducing the bias.

\subsection{Recurrent Rollout Buffer and Sequence Replay}
\label{sec:rl-buffer}

That the network carries memory between observations --- which is the reason it
was chosen --- has consequences for the organisation of the collected data and
for the update, and it is what distinguishes this realisation from a stateless
implementation.

\paragraph{The unit of a mini-batch is a sequence, not a sample.}
A mini-batch counts sequences of consecutive steps of a single environment; only
the stateless path shuffles individual samples. The buffer stores, for every
step and besides the transition, the snapshot of the recurrent state taken
\emph{before} the forward pass, since it is the state that produced that action.

\paragraph{Backpropagation is truncated at the window.}
The update replays every sequence forwards, one step at a time, feeding the
network its own output state as the next input: the gradient flows through the
length of the window and no further. The initial condition is the snapshot
stored at the start of the window.

\paragraph{The stored action is the sampled one.}
The environment clips the command within the limits declared by the platform
before applying it, but the buffer keeps the action as the policy sampled it.
The stored log-probability must remain that of \emph{that} action, or the ratio
of Equation~\eqref{eq:ppo-ratio} would be a ratio between two different
quantities.

\paragraph{What is not real must not contribute.}
The last window of a rollout is usually shorter than the nominal length, and
every loss is therefore masked and divided by the number of real positions
rather than by the size of the tensor. Likewise, an episode boundary inside a
window clears the state of the columns that ended and of those alone, so that
the memory of one episode does not leak into the next. The batch, finally, is
never split across environments: the cells hold every environment's state
together, and splitting it would fragment that state.

\subsection{Actor, Critic and Shared Encoder}
\label{sec:rl-networks}

The actor is one of the networks of Chapter~\ref{ch:lnn}, built from the same
table the supervised training uses: there is no separate actor hierarchy, since
the base class of the networks already provides the handling of the recurrent
state on which the algorithm relies.

The critic is instead a network of its own and deliberately distinct from the
actor. On an environment that includes the point cloud it also receives the
descriptor of the encoder, concatenated with the dense observation: a value
function that cannot see the obstacles is estimating a different problem from
the one the policy is solving.

One point deserves attention because violating it produces numbers rather than
errors. The critic reads the observation \emph{through the actor's
normalisation} --- a view onto the latter's statistics, neither a second
normaliser nor a copy of those statistics. With a pretrained network, which
carries its own normalisation inside the weights
(Section~\ref{sec:rl-finetuning}), a critic reading the raw vector would receive
quantities orders of magnitude out of scale --- among them the model identifier
--- and would saturate at the first nonlinearity, reducing itself to predicting
a constant. The view is not registered as a submodule, so that the statistics
exist once only and a later load onto the actor cannot leave the two in
disagreement, and it takes the leading components of the observation, since the
fixed-step variants append the elapsed interval to it
(Section~\ref{sec:cell-time}).

\paragraph{The encoder: frozen or trained.}
The sparse encoder may be kept fixed or fine-tuned along with the rest. Freezing it
allows the descriptor rather than the cloud to be stored in the buffer, and the
convolution not to be run during the update. Training it costs an additional
convolution over every sample of every epoch, and the memory to hold its graph,
which is proportional to the product of sequence length and batch size and to
the volume of the grid rather than to the number of occupied voxels.

The choice is not free in one case, however: starting from random weights, a
frozen encoder produces nearly identical descriptors for any cloud, so that the
policy would fly on odometry alone and no amount of steps would correct it.
Starting from scratch and training the encoder are therefore a single choice,
and their contrary combination is refused rather than executed.

\subsection{Action Distribution and Bound Penalty}
\label{sec:rl-actions}

The action space is continuous and the policy is Gaussian: the network produces
the mean, while the spread is a trained parameter that lives on the \emph{agent}
and not on the network. The placement is motivated: without it there would be no
sampling, no log-probability, no ratio and no entropy --- it is part of the
algorithm and not of the architecture, and a network exported to the robot need
not carry it.

The spread is initialised at a fraction of the action envelope declared by the
environment, and not at a fixed value: a unit standard deviation would be wider
than the entire action space on a platform with tight limits, and on a
pretrained network it would drown the imitated policy from the first step.

A penalty term finally charges the part of the Gaussian mean that falls outside
the envelope. Without it the network may push the mean arbitrarily beyond the
limit, since the environment clips the command in any case: the policy would pay
in entropy and in exploration for a region of the action space it cannot reach.

\subsection{Vectorised Fleet Environment and Domain Randomisation}
\label{sec:rl-environment}

The environment is not a simulator but a \emph{fleet} of isolated copies of the
entire system. Each parallel environment --- a \emph{slice} --- is a complete
vertical stack: the simulator with its own world, the bridge towards the ROS
channels, and a reduced version of the pipeline of
Chapter~\ref{ch:pipeline}. Isolation between slices is obtained on three
distinct planes --- middleware domain, simulator partition, communication port
--- and the training attaches to them through one socket per slice. A property
important to the interpretation of the results follows: the observation on which
the policy is fine-tuned is produced by the same chain that produces it on board,
and not by a simplified model of the environment.

The fleet is moreover the way domain randomisation reaches this network. What
must necessarily agree across slices is whatever sizes the actor and the buffer:
the width of the dense observation, the presence of the cloud, the width of the
action, and the names of the signals the reward is built from. What must
\emph{not} agree is everything else --- the extent and the resolution of the
grid, the robot model, the goal, the kinematic limits --- and it is precisely of
those differences that the randomisation consists. The circumstance is so
structural that the voxel pitch is one of the components of the dense
observation (Section~\ref{sec:observation-vector}): the network is not required
to infer the resolution it is looking at, it is told.

The effective heterogeneity is, however, a property of the fleet manifest and
not of the code: slices are assigned cyclically to the available configurations,
and a manifest built from a single configuration produces identical slices
without anything downstream complaining. It is a condition the experimental
protocol must verify before a long session rather than infer afterwards.

\subsection{Episode Policy: Spawn, Termination and Truncation}
\label{sec:rl-episodes}

What an episode on a slice is, is decided by the environment and not by the
training, and it is entirely \emph{derived from what that slice declares}: the
maximum number of steps from the ratio between the distance of the goal and the
envelope speed, multiplied by a budget factor; the arrival tolerance as a
multiple of the voxel pitch of that slice, so that the threshold remains
commensurate with the resolution at which that slice measures; the tipping angle
as pure geometry; the seed from the slice's own episode counter.

\paragraph{The budget factor is a demand on speed.}
Since the maximum number of steps is sized on the envelope speed, allowing a
multiple of the ideal time is the same as demanding an average speed equal to
the envelope divided by that multiple. It is a claim about the policy and not
about the world, and it is to be checked against the speed the policy actually
holds: a budget the demonstrator itself would not meet cannot be met by a policy
fine-tuned from it, and it makes arrival structurally unreachable.

\paragraph{Every episode restarts from the spawn.}
An episode resuming half way along the route would be a shorter task scored on
the same scale, and the return would report an improvement where the vehicle had
merely been asked to do less. Restarting from the origin moreover preserves the
meaning of the deadline, which measures the trip from the start, and prevents a
robot stuck against an obstacle from being truncated in place indefinitely,
which would make the distribution of the initial states a function of the
policy's own failures.

\paragraph{The task belongs to the slice.}
This is a distinct question from the previous one. The spawn and the goal are
part of the configuration of each slice, like its robot and its world, and are
not drawn by the training: every slice poses the same trip at every episode.
What varies is the trip from one slice to the next, so that a fleet can pose
tasks of graded difficulty side by side --- short trips with the goal in view
next to longer ones with obstacles on the straight route --- and every batch of
data then holds them in fixed proportions. This is not a curriculum: no task
follows another, and the training does not know which slice is easy. The
fleet used in the experiments is described in Section~\ref{sec:res-rl-fleet}.

\paragraph{The terminal conditions.}
An episode ends on arrival within the tolerance, on tipping beyond the declared
angle, on contact when a collision margin is armed, or on stalling. The last
condition deserves a motivation, because it is not a geometric threshold: the
distance from obstacles is measured from the centre of the vehicle, and is
therefore in fact a threshold on the sphere enclosing its body --- an acceptable
model for a multirotor and an inadequate one for an elongated ground vehicle,
whose sphere is inside a wall whenever it merely drives past one. The
consequence of contact, however, requires no geometry: a robot commanded forward
that does not move forward has hit something. Stalling ends the episode after a
number of consecutive steps below a fraction of the commanded speed, and only
when the command exceeds a minimum threshold --- standing still is a legitimate
thing to be told to do. It is armed per slice and only where it separates the
two cases, that is on the platforms that do not fly.

Reaching the step limit finally produces a \emph{truncation}, which the
algorithm treats differently from a termination (Section~\ref{sec:rl-gae}).

\subsection{Reward Design}
\label{sec:rl-reward}

The reward is not fixed: it is a policy chosen among several, and its structure
is a list of cost terms plus a single term that earns. Every term knows its own
name, its share of the budget and its own arithmetic, and is logged separately,
so that the contribution of each is readable.

\paragraph{The term that earns.}
Progress towards the goal is formulated as potential-based shaping
\cite{ng1999policy}, with potential equal to the negative of the distance to the
goal. The canonical form, $d_{\text{prev}} - \gamma d$, is invariant with
respect to the optimal policy, but it pays $(1-\gamma)d$ to a stationary vehicle
--- \emph{the more so the further away it is}. If the discount factor is derived
from the episode length, as described in Section~\ref{sec:rl-presets}, that rate
works out at exactly the ceiling below which the dense costs must be held: no
allocation of the shares can then make loitering unprofitable, since the matter
is excluded by construction and not by tuning. The policy adopted therefore pins
at one the discount of the shaping term alone, which thus pays the metres
actually closed. What is given up is exact invariance under a discounted return;
what is obtained is that same invariant form minus a cost proportional to the
distance still to go --- loitering far costs more than loitering close, which is
a term one would write on purpose.

\paragraph{The barriers.}
Three of the ways the vehicle is lost --- contact, tipping over, leaving the
altitude band --- have the same shape: a quantity approaching a threshold past
which the episode ends. They therefore share a single term, quadratic, null up
to a soft threshold and clipped at one beyond the hard one. Each of the three
choices answers a failure of the obvious alternative. Without the soft threshold
the cost would jump from nothing to the terminal penalty, leaving no gradient
along which to learn not to arrive there. A linear ramp would tax the useful
manoeuvre as much as the dangerous one --- a multirotor tilts in order to
accelerate and passes close to a wall in order to get past it --- whereas the
quadratic leaves the useful zone almost free. And an unbounded barrier would
dominate every other term as soon as the violation is large, and would have no
finite worst case, so that neither of the two budget constraints below could be
verified at all.

The safety barrier, in particular, is measured not against a constant distance
but against the room needed to stop,
\begin{equation}
  d_{\text{safe}}(v) = d_0 + \tau v + \frac{v^2}{2 a_{\text{brake}}},
  \label{eq:stopping}
\end{equation}
where $d_0$ is the radius of the sphere enclosing the vehicle plus a standoff,
and $\tau$ is the latency of perception and actuation. The altitude barrier is
defined around a cruise altitude and expressed as a fraction of it, so that it
follows the vehicle rather than the map, and it leaves out of the term the
slices to which no vertical velocity is sent.

\paragraph{The terminals, and the scale of everything else.}
A terminal term assigns the arrival bonus --- reduced in proportion to the speed
at which the goal is reached, since arriving into the goal is not arriving ---
and the penalty for losing the vehicle. Every other magnitude is expressed as a
fraction of what the slice declares: the maximum progress for the costs, the
terminal angle for the attitude ramp, the action envelope for the command
smoothness term, the map radius for the altitude band. A faster robot, a
different step or another simulator move all of them without anyone having to
remember to.

\paragraph{The two budget constraints.}
The shares must satisfy two inequalities. Per step, their sum must be below one:
above that, the best available policy is to stop paying, by standing still or by
ending the episode at once. Per episode, the sum must be below the reciprocal of
the budget factor, since a slow success over $k$ times the ideal time pays $k$
times the costs while earning the same distance; requiring the inequality to
hold for any distance yields $\sum f < 1/k$. It is a design rule rather than a
number to be retuned for every goal, and it is verified at the construction of
the policy and at the connection to the fleet, where the length of the episodes
is known.

\paragraph{The budget is necessary and not sufficient.}
Satisfying it guarantees that no allocation of the shares pays for ending an
episode early. It says nothing about how much the terms, \emph{together}, pay
for flying rather than for standing still, and a configuration may respect the
ceiling and converge to immobility all the same, because what the constraint
cannot see is the relative magnitude between the channels. A reward
configuration is therefore \emph{priced before} the session, at the operating
points at which the fleet has been measured, answering the two questions that
decide a session before it begins: whether flying pays better than creeping, and
above what probability of arriving trying pays better than not trying.

\subsection{Observation and Reward Normalisers}
\label{sec:rl-normalisers}

The agent has three optional normalisers --- of the observation, of the reward
and of the returns --- and a fourth constraint which is not optional and has no
switch, namely that the critic read the observation through whatever reads it
for the actor (Section~\ref{sec:rl-networks}).

The observation normaliser \emph{must stay off} with a pretrained network, which
carries its own normalisation inside the weights: a second normalisation on top
would feed it an input it has never seen. The combination is refused rather than
executed. The reward normaliser divides by the running standard deviation of the
discounted return and does \emph{not} subtract its mean, since subtracting it
would change the sign of the reward.

Two properties of the recursive estimation of the statistics deserve to be
stated, because both arise from characteristics of this observation. The first
is that the estimate accumulates the sum of squared deviations and not the
variance: the two coincide only as long as every column varies, whereas in this
observation a substantial part of the columns is constant by construction ---
the goal block, the declared envelope of the robot, the model identifier --- and
carrying a variance would mean treating as a measurement the value of
convenience assigned to a constant column, then diluting it at every batch. The
second is that every sample is bounded to a number of deviations before being
folded in: there being no forgetting factor, a single reading coming from a
diverged estimator would become, for good, the scale by which everything else is
divided.

\subsection{Trust Region, Entropy and Schedules}
\label{sec:rl-schedules}

After every mini-batch the approximate divergence between the policy that
collected the data and the current one is measured, with the low-variance
estimator $\mathbb{E}[(r-1) - \log r]$. When the running mean of the epoch
exceeds a threshold, the actor stops for the rest of the update: the data no
longer describe the current policy, and the following epochs would optimise with
respect to a distribution that has moved. The critic continues, having no ratio
and nothing to be wrong about.

\paragraph{The mean, and not the sample that produced it.}
The criterion is read on the mean and not on the individual mini-batch, and this
is a substantive choice. A mini-batch is here the product of sequence length and
batch size, and on so few samples the estimator is heavy-tailed: stopping at the
first sample beyond the threshold means stopping at the maximum of many noisy
draws, which exceeds any threshold close to the mean. The region is therefore
not consulted until a minimum number of informative estimates is available,
since the standard error falls as the square root of their number. Moreover,
only the mini-batches scored \emph{after} the policy has moved enter that mean:
until the first optimiser step the weights are still the ones that collected the
rollout, the ratio is exactly one and the estimate exactly zero, and including
them would lower the mean by an amount that depends on the size of the gradient
buffer rather than on the policy.

\paragraph{The schedules.}
The learning rate and the clipping width $\epsilon$ both decay linearly towards
a fraction of their initial value rather than towards zero: early on, large and
exploratory steps are wanted; late on, a policy that settles. The horizon of the
decay is derived from the step budget, so that raising it does not merely
lengthen the session but also re-paces its schedule.

\paragraph{One knob wearing four names.}
The initial exploration width, the actor's learning rate, the trust-region
threshold and the number of mini-batches accumulated per step are not
independent parameters. The divergence grows as $\Delta\mu^2/2\sigma^2$ and the
gradient on the mean as $A(a-\mu)/\sigma^2$: the trust region is therefore
\emph{measured in units of the exploration width}, and narrowing the latter
makes the same policy step far larger in the eyes of the stopping criterion.
They must therefore be tuned together, which is why they are the only parameters
of the algorithm the command line exposes beside those of the environment.

\subsection{Fine-Tuning from Imitated Weights}
\label{sec:rl-finetuning}

Loading an imitation checkpoint rebuilds the network from the declaration held
in the metadata, by the same route the inference node takes on board: training
and inference load the same file in the same way. Three adjustments are applied
on the way in, and each of them guards against a failure that would produce
numbers rather than an exception.

\begin{enumerate}
  \item \textbf{The external normalisation stays off}, for the reason given in
        Section~\ref{sec:rl-normalisers}.
  \item \textbf{Degenerate statistics are repaired}, as in
        Section~\ref{sec:sl-normaliser}. The repair is neutral on constant data,
        and fine-tuning is precisely what makes those features start to vary.
  \item \textbf{The output calibration becomes trainable.} Kept fixed, a
        component the demonstration never exercised has zero correlation, hence
        zero gain: under fine-tuning that factor multiplies the gradient too, so
        that the component is not attenuated but \emph{disconnected} --- and
        those are exactly the components pretraining got wrong. The trainable
        mode starts from the estimated values and lets them move.
\end{enumerate}

\paragraph{The first iterations.}
A pretrained actor arrives with a policy worth keeping and a critic that knows
nothing. The first advantages are therefore the critic's error, and the
algorithm would spend the first iterations moving good weights with them --- the
encoder included, where it is trainable. A configurable number of initial
iterations is therefore reserved for the critic alone: while it lasts, the
actor's weights, its exploration width and the auxiliary head remain unchanged
bit for bit. The phase has moreover an experimental value of its own, and it is
stated here because Chapter~\ref{ch:results} relies on it: for its whole
duration the performance metrics describe the imitated policy exactly as
imitation left it, and constitute therefore its only free measurement on the
fleet.

\paragraph{Starting over or continuing.}
Starting from the weights alone begins a new session, with critic, optimisers
and exploration width afresh; continuing an interrupted one also restores them,
together with the statistics of the normalisers and the position in the
schedule, each of which would otherwise make the resumption read as a regression
that did not occur.

\paragraph{What can be transferred on board.}
A policy fine-tuned from imitation carries the statistics of the corpus inside its
own weights and is transferable as it stands. A policy trained from random
weights is not: its internal normalisation is the identity, and the
normalisation that made the training work lives in the agent's recursive
estimator, which inference never opens. In the second case the statistics must
be written into the network before export, an operation the dedicated tool
performs while refusing the cases in which it would be ambiguous --- a
checkpoint carrying them in both places, or in neither.

\subsection{Hyperparameter Presets}
\label{sec:rl-presets}

The hyper-parameters of the algorithm are declared as data, one set per
architecture, rather than as arguments scattered among the call sites: changing
architecture is a choice on the command line, changing a value is an edit in a
single visible place. Each set moreover provides a \emph{substitution} for
starting from random weights, which alters the few values that regime requires
to differ --- a higher learning rate and entropy bonus, a wider initial
exploration, no phase reserved for the critic, there being no policy to protect,
and the external normalisation of the observation switched on, since in that
regime the network carries none of its own inside. It is the same difference
that Section~\ref{sec:rl-finetuning} states with respect to transfer on board.
The command line reaches only the parameters on which a fleet is genuinely
retuned, and every value differing from the default set is printed at startup,
so that the session itself declares what it was started with.

\paragraph{One discount per slice.}
The discount factor deserves separate treatment, because a heterogeneous fleet
admits no single one. Resolving it as $\gamma = 1 - 1/T$ from the episode length
$T$ of each slice makes $\gamma^{T}$ equal to $1/e$ on all of them: the arrival
bonus is equally visible from every spawn. A single value derived from the mean
would instead give slices with very different trips very different horizons,
posing in effect the same task as two distinct problems, and making the bonus
nearly invisible to the longest slice. The same per-slice value is read by the
advantage trace, by the value added on truncation, by the reward normaliser and
by the progress potential, so that nothing is discounted twice with two
different values.
